\documentclass{article}
\usepackage[utf8]{inputenc}
\usepackage[T2A]{fontenc}
\usepackage[russian]{babel}
\usepackage{indentfirst}
\usepackage{amsmath}
\usepackage{amssymb}

\sloppy

\begin{document}

\section*{Векторное исчисление}

\subsection*{Упражнение 5.5}

\subsubsection*{Условие}

Рассмотрим следующие функции:
\begin{align*}
f_1(x) &= sin(x_1) cos(x_2), \quad x \in \mathbb{R}^2 \\
f_2(x, y) &= x^T y, \quad x, y \in \mathbb{R}^n \\
f_3(x) &= x x^T, \quad x \in \mathbb{R}^n
\end{align*}

а. Каковы размерности матрицы \(\frac{\partial f_i}{\partial x}\)?

б. Вычислить матрицы Якоби.

\subsubsection*{Решение}

1. Функция \( f_1(x) = sin(x_1) cos(x_2), \quad x \in \mathbb{R}^2 \).

\quad

Частные производные:
\begin{align*}
\frac{\partial f_1}{\partial x_1} &= \cos(x_1)\cos(x_2) \\
\frac{\partial f_1}{\partial x_2} &= -\sin(x_1)\sin(x_2)
\end{align*}

Матрица Якоби:
\[
J_{f_1} = \frac{\partial f_1}{\partial x} = \begin{bmatrix}
\dfrac{\partial f_1}{\partial x_1} & \dfrac{\partial f_1}{\partial x_2}
\end{bmatrix} = \begin{bmatrix}
\cos(x_1)\cos(x_2) & -\sin(x_1)\sin(x_2)
\end{bmatrix}
\]

Таким образом, получили матрицу \(\frac{\partial f_1}{\partial x}\) размерностью \(1 \times 2\).

\quad

2. Функция \( f_2(x, y) = x^T y, \quad x, y \in \mathbb{R}^n \).

\quad

Координатная форма:
\[
f_2(x, y) = x^T y = \sum_{i=1}^n x_i y_i
\]

Частные производные:
\begin{align*}
\frac{\partial f_2}{\partial x_i} &= y_i \quad \text{для } i = 1, \ldots, n \\
\frac{\partial f_2}{\partial y_j} &= x_j \quad \text{для } j = 1, \ldots, n
\end{align*}

Матрица Якоби:
\[
J_{f_2} = \frac{\partial f_2}{\partial (x, y)} = \begin{bmatrix}
\dfrac{\partial f_2}{\partial x} & \dfrac{\partial f_2}{\partial y}
\end{bmatrix} = \begin{bmatrix} y^T & x^T \end{bmatrix}
\]

В развёрнутой форме:
\[
J_{f_2} = \begin{bmatrix}
y_1 & y_2 & \cdots & y_n & x_1 & x_2 & \cdots & x_n
\end{bmatrix}
\]

Таким образом, получили матрицу \(\frac{\partial f_2}{\partial x}\) размерностью \(1 \times 2n\).

\quad

3. Функция \( f_3(x) = x x^T, \quad x \in \mathbb{R}^n \).

\quad

Запись в матричном виде:
\[
f_3(x) = x x^T = \begin{bmatrix}
x_1^2 & x_1x_2 & \cdots & x_1x_n \\
x_2x_1 & x_2^2 & \cdots & x_2x_n \\
\vdots & \vdots & \ddots & \vdots \\
x_nx_1 & x_nx_2 & \cdots & x_n^2
\end{bmatrix}
\]

Рассмотрим произвольный элемент матрицы: \([f_3(x)]_{ij} = x_i x_j\).

\quad

Вычислим частные производные этого элемента по всем переменным \(x_k\): 
\begin{align*}
\frac{\partial (x_i x_j)}{\partial x_i} &= x_j, \quad \text{при } k = i, k \neq j \\
\frac{\partial (x_i x_j)}{\partial x_j} &= x_i, \quad \text{при } k = j, k \neq i \\
\frac{\partial (x_i^2)}{\partial x_i} &= 2x_i, \quad \text{при } k = i = j \\
\frac{\partial (x_i x_j)}{\partial x_k} &= 0, \quad \text{при } k \neq i, k \neq j
\end{align*}

Пусть \(x = \begin{bmatrix} x_1 \\ x_2 \end{bmatrix}\), тогда:
\[
f_3(x) = x x^T = \begin{bmatrix}
x_1^2 & x_1x_2 \\
x_1x_2 & x_2^2
\end{bmatrix}
\]

Вычислим матрицу Якоби в этом случае. Для этого сначала запишем все элементы функции в виде вектора:
\[
\text{vec}(f_3(x)) = \begin{bmatrix}
x_1^2 \\ x_1x_2 \\ x_1x_2 \\ x_2^2
\end{bmatrix}
\]

Теперь вычислим частные производные:
\begin{align*}
\frac{\partial x_1^2}{\partial x_1} &= 2x_1 &
\frac{\partial x_1^2}{\partial x_2} &= 0 \\
\frac{\partial (x_1x_2)}{\partial x_1} &= x_2 &
\frac{\partial (x_1x_2)}{\partial x_2} &= x_1 \\
\frac{\partial (x_1x_2)}{\partial x_1} &= x_2 &
\frac{\partial (x_1x_2)}{\partial x_2} &= x_1 \\
\frac{\partial x_2^2}{\partial x_1} &= 0 &
\frac{\partial x_2^2}{\partial x_2} &= 2x_2
\end{align*}

Таким образом, получаем матрицу Якоби для \(n = 2\):
\[
J_{f_3} = \begin{bmatrix}
2x_1 & 0 \\
x_2 & x_1 \\
x_2 & x_1 \\
0 & 2x_2
\end{bmatrix}
\]

В общем случае матрица Якоби имеет размерность \(n^2 \times n\). Её можно построить по блокам. Для каждого \(k = 1, \ldots, n\), столбец матрицы Якоби, соответствующий \(\frac{\partial f_3}{\partial x_k}\), представляет собой матрицу \(n \times n\):

\[
\frac{\partial f_3}{\partial x_k} = e_k x^T + x e_k^T
\]
где \(e_k\) — \(k\)-й базисный вектор.

\quad

В развернутой форме получаем следующее:
\[
\frac{\partial f_3}{\partial x_k} = \begin{bmatrix}
0 & \cdots & 0 & x_1 & 0 & \cdots & 0 \\
\vdots & \ddots & \vdots & \vdots & \vdots & \ddots & \vdots \\
0 & \cdots & 0 & x_{k-1} & 0 & \cdots & 0 \\
x_1 & \cdots & x_{k-1} & 2x_k & x_{k+1} & \cdots & x_n \\
0 & \cdots & 0 & x_{k+1} & 0 & \cdots & 0 \\
\vdots & \ddots & \vdots & \vdots & \vdots & \ddots & \vdots \\
0 & \cdots & 0 & x_n & 0 & \cdots & 0
\end{bmatrix}
\]

\subsection*{Упражнение 5.6}

\subsubsection*{Условие}

Дифференцируйте \(f\) относительно \(t\) и \(g\) относительно \(X\), где
\begin{align*}
f(t) &= \sin(\log(t^T t)), \quad t \in \mathbb{R}^D \\
g(X) &= \operatorname{tr}(AXB), \quad A \in \mathbb{R}^{D \times E}, X \in \mathbb{R}^{E \times F}, B \in \mathbb{R}^{F \times D}
\end{align*}

\subsubsection*{Решение}

1. Функция \( f(t) = \sin(\log(t^T t)), \quad t \in \mathbb{R}^D \).

\quad

Запишем её как композицию:
\[
f(t) = \sin(u(t)), \quad \text{где } u(t) = \log(v(t)), \quad \text{и } v(t) = t^T t
\]

Вычислим частные производные каждой функции:
\begin{align*}
\frac{\partial v}{\partial t} &= \frac{\partial (t^T t)}{\partial t} = 2t \\
\frac{\partial u}{\partial v} &= \frac{\partial \log(v)}{\partial v} = \frac{1}{v} = \frac{1}{t^T t} \\
\frac{\partial f}{\partial u} &= \frac{\partial \sin(u)}{\partial u} = \cos(u) = \cos(\log(t^T t))
\end{align*}

И применим правило цепочки:
\[
\frac{\partial f}{\partial t} = \frac{\partial f}{\partial u} \cdot \frac{\partial u}{\partial v} \cdot \frac{\partial v}{\partial t} = \cos(\log(t^T t)) \cdot \frac{1}{t^T t} \cdot 2t = \frac{2\cos(\log(t^T t))}{t^T}
\]

Таким образом:
\[
\frac{\partial f}{\partial t} = \frac{2\cos(\log(t^T t))}{t^T}
\]

\quad

2. Функция \( g(X) = \operatorname{tr}(AXB) \), где \( A \in \mathbb{R}^{D \times E}, X \in \mathbb{R}^{E \times F}, B \in \mathbb{R}^{F \times D} \).

\quad

Запишем след в координатной форме:
\[
\operatorname{tr}(AXB) = \sum_{i=1}^D \sum_{j=1}^E \sum_{k=1}^F A_{ij} X_{jk} B_{ki}
\]

Вычислим частную производную по элементу \( X_{mn} \):
\[
\frac{\partial g}{\partial X_{mn}} = \sum_{i=1}^D \sum_{j=1}^E \sum_{k=1}^F A_{ij} \frac{\partial X_{jk}}{\partial X_{mn}} B_{ki}
\]

Здесь мы можем воспользоваться свойством \( \frac{\partial X_{jk}}{\partial X_{mn}} = \delta_{jm} \delta_{kn} \), где \( \delta_{ij} \) — символ Кронекера, который определяется как:
\[
\delta_{ij} = \begin{cases}
1 & \text{если } i = j \\
0 & \text{если } i \neq j
\end{cases}
\]

Благодаря этому свойству двойная сумма упрощается, потому что \( \frac{\partial X_{jk}}{\partial X_{mn}} \neq 0 \) только тогда, когда \( j = m \) и \( k = n \).

Следовательно, получаем:
\[
\frac{\partial g}{\partial X_{mn}} = \sum_{i=1}^D A_{im} B_{ni}
\]

Это соответствует элементу матрицы \( B^T A^T \), но поскольку мы хотим получить \( \frac{\partial g}{\partial X} \), а не \( \frac{\partial g}{\partial X^T} \), имеем:
\[
\frac{\partial g}{\partial X} = A^T B^T
\]

\subsection*{Упражнение 5.7}

\subsubsection*{Условие}

Вычислите производные \( \frac{\partial f}{\partial x} \) следующих функций, используя правило цепочки. Укажите размеры каждой отдельной частной производной. Подробно опишите свои действия.

\quad

а.
\[
f(z) = \log(1 + z), \quad z = x^T x, \quad x \in \mathbb{R}^D
\]

\quad

б.
\[
f(z) = \sin(z), \quad z = Ax + b, \quad A \in \mathbb{R}^{E \times D}, x \in \mathbb{R}^D, b \in \mathbb{R}^E
\]
где \(\sin(\cdot)\) применяется к каждому элементу \(z\).

\subsubsection*{Решение}

а. Функция \( f(z) = \log(1 + z), \quad z = x^T x, \quad x \in \mathbb{R}^D \).

\quad

Функция представляет собой композицию:
\[
f(x) = h(g(x)), \quad \text{где } h(z) = \log(1 + z), \quad g(x) = x^T x
\]

Вычислим промежуточные частные производные.

\quad

Производная \( \frac{\partial h}{\partial z} \):
\[
\frac{\partial h}{\partial z} = \frac{\partial}{\partial z} \log(1 + z) = \frac{1}{1 + z} \quad (1 \times 1)
\]

Производная \( \frac{\partial z}{\partial x} \):
\[
z = x^T x = \sum_{i=1}^D x_i^2
\]
\[
\frac{\partial z}{\partial x_j} = \frac{\partial}{\partial x_j} \left( \sum_{i=1}^D x_i^2 \right) = 2x_j
\]
\[
\frac{\partial z}{\partial x} = 2x \quad (D \times 1)
\]

Применим правило цепочки:
\[
\frac{df}{dx} = \frac{\partial h}{\partial z} \cdot \frac{\partial z}{\partial x}
= \frac{1}{1 + z} \cdot 2x
= \frac{2x}{1 + x^T x}
\]

\quad

б. Функция \( f(z) = \sin(z), \quad z = Ax + b, \quad A \in \mathbb{R}^{E \times D}, x \in \mathbb{R}^D, b \in \mathbb{R}^E \).

\quad

Функция представляет собой композицию:
\[
f(x) = h(g(x)), \quad \text{где } h(z) = \sin(z), \quad g(x) = Ax + b
\]

Вычислим промежуточные частные производные.

\quad

Поскольку \(\sin(\cdot)\) применяется поэлементно, производная \( \frac{\partial h}{\partial z} \) представляет собой диагональную матрицу:
\[
\frac{\partial h}{\partial z} = \text{diag}(\cos(z)) = \begin{bmatrix}
\cos(z_1) & 0 & \cdots & 0 \\
0 & \cos(z_2) & \cdots & 0 \\
\vdots & \vdots & \ddots & \vdots \\
0 & 0 & \cdots & \cos(z_E)
\end{bmatrix} \quad (E \times E)
\]

Производная \( \frac{\partial z}{\partial x} \):
\[
z = Ax + b \Rightarrow \frac{\partial z}{\partial x} = A \quad (E \times D)
\]

Применим правило цепочки:
\[
\frac{df}{dx} = \frac{\partial h}{\partial z} \cdot \frac{\partial z}{\partial x}
= \text{diag}(\cos(z)) \cdot A
\]

\subsection*{Упражнение 5.8}

\subsubsection*{Условие}

Вычислите производные \( \frac{\partial f}{\partial x} \) следующих функций. Подробно опишите свои действия.

\quad

а. Используйте правило цепочки. Укажите размеры каждой отдельной частной производной.
\begin{align*}
f(z) &= \exp\left(-\frac{1}{2}z\right) \\
z &= g(y) = y^T S^{-1} y \\
y &= h(x) = x - \mu
\end{align*}
где \( x, \mu \in \mathbb{R}^D, S \in \mathbb{R}^{D \times D} \).

\quad

б.
\[
f(x) = \operatorname{tr}(xx^T + \sigma^2 I), \quad x \in \mathbb{R}^D
\]

\quad

в. Используйте правило цепочки. Укажите размеры каждой отдельной частной производной. Вам не нужно вычислять произведение частных производных явно.
\begin{align*}
f &= \tanh(z) \in \mathbb{R}^M \\
z &= Ax + b, \quad x \in \mathbb{R}^N, A \in \mathbb{R}^{M \times N}, b \in \mathbb{R}^M
\end{align*}
где \(\tanh\) применяется поэлементно к каждому компоненту \(z\).

\subsubsection*{Решение}

a. Определим структуру композиции:
\[
f(x) = f(z(g(y(h(x)))))
\]

Вычислим все частные производные.

\quad

Производная \( \frac{\partial f}{\partial z} \):
\[
\frac{\partial f}{\partial z} = \frac{\partial}{\partial z} \exp\left(-\frac{1}{2}z\right) = -\frac{1}{2} \exp\left(-\frac{1}{2}z\right) \quad (1 \times 1)
\]

Производная \( \frac{\partial z}{\partial y} \):
\[
z = y^\top S^{-1} y = \sum_{i=1}^D \sum_{j=1}^D y_i [S^{-1}]_{ij} y_j
\]
\[
\frac{\partial z}{\partial y_k} = \sum_{j=1}^D [S^{-1}]_{kj} y_j + \sum_{i=1}^D y_i [S^{-1}]_{ik} = 2 \sum_{j=1}^D [S^{-1}]_{kj} y_j
\]
\[
\frac{\partial z}{\partial y} = 2 S^{-1} y \quad (1 \times D)
\]

Производная \( \frac{\partial y}{\partial x} \):
\[
y = x - \mu \Rightarrow \frac{\partial y}{\partial x} = I_D \quad (D \times D)
\]

Применим правило цепочки:
\[
\frac{df}{dx} = \frac{\partial f}{\partial z} \cdot \frac{\partial z}{\partial y} \cdot \frac{\partial y}{\partial x}
\]
\[
\frac{df}{dx} = \left(-\frac{1}{2} \exp\left(-\frac{1}{2}z\right)\right) \cdot (2 S^{-1} y) \cdot I_D
\]
\[
\frac{df}{dx} = -\exp\left(-\frac{1}{2}y^T S^{-1} y\right) \cdot S^{-1} y
\]

Таким образом, получаем что:
\[
\frac{df}{dx} = -\exp\left(-\frac{1}{2}(x - \mu)^T S^{-1} (x - \mu)\right) \cdot S^{-1} (x - \mu)
\]

\quad

б. Функция \( f(x) = \operatorname{tr}(xx^T + \sigma^2 I), \quad x \in \mathbb{R}^D \).

\quad

Запишем внешнее произведение явно:
\[
x x^T = \begin{bmatrix}
x_1^2 & x_1 x_2 & \cdots & x_1 x_D \\
x_2 x_1 & x_2^2 & \cdots & x_2 x_D \\
\vdots & \vdots & \ddots & \vdots \\
x_D x_1 & x_D x_2 & \cdots & x_D^2
\end{bmatrix}
\]

Вычислим след:
\[
f(x) = \operatorname{tr}(x x^T + \sigma^2 I) = \operatorname{tr}(x x^T) + \operatorname{tr}(\sigma^2 I)
\]
\[
\operatorname{tr}(x x^T) = \sum_{i=1}^D x_i^2 = x^T x
\]
\[
\operatorname{tr}(\sigma^2 I) = \sigma^2 D
\]
\[
f(x) = x^T x + \sigma^2 D
\]

Вычислим производную:
\[
\frac{df}{dx} = \frac{\partial}{\partial x} (x^T x + \sigma^2 D) = 2x
\]

\quad

в. Определим структуру композиции:
\[
f(x) = \tanh(Ax + b)
\]

Вычислим частные производные.

\quad

Поскольку \(\tanh\) применяется поэлементно, производная \( \frac{\partial f}{\partial z} \) представляет собой диагональную матрицу:
\[
\frac{\partial f}{\partial z} = \text{diag}\left(1 - \tanh^2(z)\right) \quad (M \times M)
\]
\[
\frac{\partial f}{\partial z} = \begin{bmatrix}
1 - \tanh^2(z_1) & 0 & \cdots & 0 \\
0 & 1 - \tanh^2(z_2) & \cdots & 0 \\
\vdots & \vdots & \ddots & \vdots \\
0 & 0 & \cdots & 1 - \tanh^2(z_M)
\end{bmatrix}
\]

Производная \( \frac{\partial z}{\partial x} \):
\[
z = Ax + b \Rightarrow \frac{\partial z}{\partial x} = A \quad (M \times N)
\]

Применим правило цепочки:
\[
\frac{df}{dx} = \frac{\partial f}{\partial z} \cdot \frac{\partial z}{\partial x} = \text{diag}\left(1 - \tanh^2(Ax + b)\right) \cdot A
\]

\subsection*{Упражнение 5.9}

\subsubsection*{Условие}

Определим:
\begin{align*}
g(x, z, \nu) &:= \log p(x, z) - \log q(z, \nu) \\
z &:= t(\epsilon, \nu)
\end{align*}
где \(p, q, t\) — дифференцируемые функции, и \(x \in \mathbb{R}^D\), \(z \in \mathbb{R}^E\), \(\nu \in \mathbb{R}^F\), \(\epsilon \in \mathbb{R}^G\).

\quad

Требуется вычислить следующий градиент, используя правило цепочки:
\[
\frac{d}{d\nu} g(x, z, \nu)
\]

\subsubsection*{Решение}

Запишем полную производную, используя правило цепочки:
\[
\frac{dg}{d\nu} = \frac{\partial g}{\partial \nu} + \frac{\partial g}{\partial z} \cdot \frac{\partial z}{\partial \nu}
\]

Вычислим частные производные.

\quad

Производная \(\frac{\partial g}{\partial \nu}\):
\begin{align*}
\frac{\partial g}{\partial \nu} &= \frac{\partial}{\partial \nu} (\log p(x, z) - \log q(z, \nu)) \\
&= 0 - \frac{\partial}{\partial \nu} \log q(z, \nu) \\
&= -\frac{1}{q(z, \nu)} \cdot \frac{\partial q(z, \nu)}{\partial \nu} \quad (1 \times F)
\end{align*}

Производная \(\frac{\partial g}{\partial z}\):
\begin{align*}
\frac{\partial g}{\partial z} &= \frac{\partial}{\partial z} (\log p(x, z) - \log q(z, \nu)) \\
&= \frac{1}{p(x, z)} \cdot \frac{\partial p(x, z)}{\partial z} - \frac{1}{q(z, \nu)} \cdot \frac{\partial q(z, \nu)}{\partial z} \quad (1 \times E)
\end{align*}

Производная \(\frac{\partial z}{\partial \nu}\):
\[
\frac{\partial z}{\partial \nu} = \frac{\partial t(\epsilon, \nu)}{\partial \nu} \quad (E \times F)
\]

Подставлим в формулу полной производной:
\[
\frac{dg}{d\nu} = -\frac{1}{q(z, \nu)} \cdot \frac{\partial q(z, \nu)}{\partial \nu} + (\frac{1}{p(x, z)} \cdot \frac{\partial p(x, z)}{\partial z} - \frac{1}{q(z, \nu)} \cdot \frac{\partial q(z, \nu)}{\partial z}) \cdot \frac{\partial t(\epsilon, \nu)}{\partial \nu}
\]

\end{document}